\documentclass[aspectratio=169]{beamer}
\usetheme{Madrid}
\usecolortheme{default}

% Packages
\usepackage{amsmath,amssymb}
\usepackage{graphicx}
\usepackage{booktabs}
\usepackage{subcaption}
\usepackage{natbib}

% Custom commands
\newcommand{\R}{\mathbb{R}}
\newcommand{\sigmoid}{\sigma}

% Title information
\title[GC24F-03]{GC24F-03 Modeling Adaptive Human Decision-Making in Refugee Movement Using Dual Process Theory (Invited)}
\subtitle{System 1 vs System 2 Processing in Forced Displacement}
\author{Michael Puma}
\institute{Columbia University Climate School\\Center for Climate Systems Research}
\date{\today}

\begin{document}

%==============================================================================
% SLIDE 1: Title
%==============================================================================
\begin{frame}
\titlepage
\end{frame}

%==============================================================================
% SLIDE 2: The Question & Dual-Process Theory
%==============================================================================
\begin{frame}{How Do People Decide Where to Go When Fleeing Conflict?}

\textbf{Two cognitive modes:}

\begin{columns}[T]
\column{0.48\textwidth}
\textbf{System 1 (S1)}
\begin{itemize}
    \item Fast, automatic
    \item Heuristic-based
    \item ``Head toward nearest border''
    \item Low cognitive load
\end{itemize}

\vspace{0.5em}
\textit{Reactive, instinctive}

\column{0.48\textwidth}
\textbf{System 2 (S2)}
\begin{itemize}
    \item Slow, deliberative
    \item Analytical
    \item ``Calculate routes, compare destinations''
    \item High cognitive load
\end{itemize}

\vspace{0.5em}
\textit{Effortful, strategic}
\end{columns}

\vspace{1em}

\textbf{Core principle: Context-dependent processing}
\begin{itemize}
    \item \textbf{High conflict} → System 1 (hard constraint)
    \item \textbf{Low conflict + social connections} → System 2 possible
\end{itemize}

\vspace{0.5em}

\centering
\alert{System 2 isn't always ``better'' – it requires mental effort that people may not have during crisis}

\end{frame}

%==============================================================================
% SLIDE 3: Our Model
%==============================================================================
\begin{frame}{Our Model: Context-Dependent Cognitive Processing}

\textbf{Key principle:} High conflict constrains people to System 1 processing, while recovery periods allow socially connected individuals to potentially engage System 2.

\vspace{0.5em}

\textbf{Mathematical framework:}
\begin{equation*}
P_{S2} = \frac{1}{1 + e^{-\eta(P(t) - \theta)}}
\end{equation*}
where $P(t) = B(t) + C(t) + S(t)$ (base + conflict + social pressure)

\vspace{0.5em}

\textbf{Key factors:}
\begin{itemize}
    \item \textbf{Conflict intensity} ($C$): High conflict → System 1 (hard constraint)
    \item \textbf{Social connections} ($S$): Connections $\geq 2$ enable System 2 capability
    \item \textbf{Recovery periods}: Lower conflict allows pressure to decrease
\end{itemize}

\vspace{0.5em}

\textbf{Implementation:} Agent-based model (FLEE framework) with four refugee scenarios:
\begin{enumerate}
    \item Nearest Border (S1 heuristic)
    \item Multiple Routes (S2 deliberation)
    \item Social Connections (S2 enabler)
    \item Context Transition (S1→S2)
\end{enumerate}

\end{frame}

%==============================================================================
% SLIDE 4: Results - Scenario Overview
%==============================================================================
\begin{frame}{Results: Four Refugee Scenarios}

\begin{figure}[htbp]
    \centering
    \begin{subfigure}[b]{0.48\textwidth}
        \centering
        \includegraphics[width=\textwidth]{run_Nearest_Border_20251211_000126_network.png}
        \caption{Nearest Border (S1)}
    \end{subfigure}
    \begin{subfigure}[b]{0.48\textwidth}
        \centering
        \includegraphics[width=\textwidth]{run_Multiple_Routes_20251211_000127_network.png}
        \caption{Multiple Routes (S2)}
    \end{subfigure}
    \vspace{0.3cm}
    \begin{subfigure}[b]{0.48\textwidth}
        \centering
        \includegraphics[width=\textwidth]{run_Social_Connections_20251211_000127_network.png}
        \caption{Social Connections}
    \end{subfigure}
    \begin{subfigure}[b]{0.48\textwidth}
        \centering
        \includegraphics[width=\textwidth]{run_Context_Transition_20251211_000128_network.png}
        \caption{Context Transition}
    \end{subfigure}
\end{figure}

\textbf{Key finding:} Context-dependent S1/S2 activation
\begin{itemize}
    \item High conflict → System 1; Recovery → System 2
    \item Social connections enable System 2 capability
    \item Same agents switch cognitive modes based on context
\end{itemize}

\end{frame}

%==============================================================================
% SLIDE 5: Results - Decision Patterns
%==============================================================================
\begin{frame}{Results: Decision Patterns}

\begin{figure}[htbp]
    \centering
    \begin{subfigure}[b]{0.48\textwidth}
        \centering
        \includegraphics[width=\textwidth]{run_Nearest_Border_20251211_000126_decisions.png}
        \caption{Nearest Border}
    \end{subfigure}
    \begin{subfigure}[b]{0.48\textwidth}
        \centering
        \includegraphics[width=\textwidth]{run_Multiple_Routes_20251211_000127_decisions.png}
        \caption{Multiple Routes}
    \end{subfigure}
    \vspace{0.3cm}
    \begin{subfigure}[b]{0.48\textwidth}
        \centering
        \includegraphics[width=\textwidth]{run_Social_Connections_20251211_000127_decisions.png}
        \caption{Social Connections}
    \end{subfigure}
    \begin{subfigure}[b]{0.48\textwidth}
        \centering
        \includegraphics[width=\textwidth]{run_Context_Transition_20251211_000128_decisions.png}
        \caption{Context Transition}
    \end{subfigure}
\end{figure}

\textbf{Key findings:}
\begin{itemize}
    \item \textbf{Route choices differ}: S1 uses heuristics (shorter routes); S2 deliberates (safer routes)
    \item \textbf{S2 activation increases} in recovery zones (low conflict + connections)
    \item \textbf{Transitions occur}: Same agents switch cognitive modes based on context
\end{itemize}

\end{frame}

%==============================================================================
% SLIDE 6: Dimensionless Parameters
%==============================================================================
\begin{frame}{Dimensionless Parameters \& Generalizability}

\textbf{Why dimensionless?} Enable comparison across scales (village vs. country-level)

\vspace{0.5em}

\begin{columns}[T]
\column{0.48\textwidth}
\begin{figure}[htbp]
    \centering
    \includegraphics[width=\textwidth]{dimensionless_parameters.png}
\end{figure}

\column{0.48\textwidth}
\textbf{Key parameters:}
\begin{itemize}
    \item $d^* = d/D_{char}$ (normalized distance)
    \item $v^* = v/v_{char}$ (normalized speed)
    \item $t^* = t/T_{char}$ (normalized time)
    \item $\eta^*$ (evacuation rate)
    \item $C^*$ (normalized conflict)
    \item $S^*$ (social connectivity)
\end{itemize}

\vspace{0.5em}

\textbf{Critical insight:}
\begin{itemize}
    \item Universal threshold: $\Theta^* \approx 0.3$ for S2 activation
    \item $C^* < 0.3$ + $S^* > 0.3$ → S2 activation
    \item Patterns hold across spatial scales
\end{itemize}
\end{columns}

\end{frame}

%==============================================================================
% SLIDE 7: Heuristic Decision-Making Patterns
%==============================================================================
\begin{frame}{Heuristic Decision-Making Patterns}

\begin{figure}[htbp]
    \centering
    \includegraphics[width=0.85\textwidth]{heuristic_decision_making.png}
\end{figure}

\textbf{Critical insights:}
\begin{itemize}
    \item \textbf{S2 activation vs. conflict}: Strong inverse relationship (high conflict → low S2)
    \item \textbf{Route choice heuristics}: S1 agents choose shorter routes; S2 agents optimize safety
    \item \textbf{Context transition}: Clear S1 → S2 transition as conflict decreases
    \item \textbf{Universal threshold}: $\Theta^* \approx 0.3$ for S2 activation (consistent across scenarios)
\end{itemize}

\end{frame}

%==============================================================================
% SLIDE 8: South Sudan Application
%==============================================================================
\begin{frame}{Application: South Sudan Displacement}

\textbf{Survey data:} ``Why do Refugees Move After Arrival? Opportunity and Community''\\
\textit{Population, Space and Place} (DOI: 10.1002/psp.2842)

\vspace{0.5em}

\textbf{Key empirical findings:}
\begin{itemize}
    \item \textbf{17\% mobility} within first year (vs. 3.4\% non-refugees)
    \item \textbf{Drivers}: Employment opportunities and established community networks
    \item \textbf{Timing}: Most moves occur \textit{after} initial settlement (not immediate flight)
\end{itemize}

\vspace{0.5em}

\textbf{S1/S2 interpretation:}
\begin{itemize}
    \item \textbf{Early phase} (active conflict): System 1 dominates → rapid, heuristic decisions
    \item \textbf{Post-arrival phase}: System 2 becomes possible → employment/community comparison
    \item \textbf{Social networks enable deliberation}: Community presence → S2 activation
\end{itemize}

\vspace{0.5em}

\textbf{``S2 is Lazy'' insight:} 17\% engage S2 when conditions right (community networks reduce cognitive cost, employment justifies effort); 83\% don't move (S2 not worth it or not possible)

\end{frame}

%==============================================================================
% SLIDE 9: Humanitarian Implications
%==============================================================================
\begin{frame}{Humanitarian Implications}

\textbf{Standard assumption:} Refugees always make calculated decisions OR always react instinctively

\vspace{0.5em}

\textbf{Our finding:} Cognitive processing varies with context

\begin{columns}[T]
\column{0.48\textwidth}
\textbf{Active conflict:}
\begin{itemize}
    \item System 1 (reactive)
    \item Simple, fast routes needed
    \item Clear, visible directions
    \item Don't assume deliberation
\end{itemize}

\column{0.48\textwidth}
\textbf{Recovery periods:}
\begin{itemize}
    \item System 2 possible (if connected)
    \item Decision support valuable
    \item Social networks enable deliberation
    \item Multiple route options allow comparison
\end{itemize}
\end{columns}

\vspace{0.5em}

\textbf{Practical recommendations:}
\begin{itemize}
    \item \textbf{Phase-based interventions}: Simple routes (active conflict) vs. decision support (recovery)
    \item \textbf{Social connection support}: Community spaces, family reunification, information sharing
    \item \textbf{Context-aware modeling}: Don't assume uniform decision-making
\end{itemize}

\end{frame}

%==============================================================================
% SLIDE 10: Key Contributions
%==============================================================================
\begin{frame}{Key Contributions}

\begin{enumerate}
    \item \textbf{Context-dependent cognitive processing}
    \begin{itemize}
        \item High conflict constrains to System 1; recovery allows System 2
        \item Same agents process differently in different contexts
    \end{itemize}
    
    \item \textbf{Social connections enable System 2}
    \begin{itemize}
        \item Not just information – connections enable mental effort
        \item Social resources shape cognitive availability
    \end{itemize}
    
    \item \textbf{Implementation in Flee}
    \begin{itemize}
        \item Agent-based model with dual-process decision-making
        \item Context-dependent rules and social connection tracking
    \end{itemize}
    
    \item \textbf{Humanitarian relevance}
    \begin{itemize}
        \item Phase-based intervention strategies
        \item Context-aware modeling approach
    \end{itemize}
\end{enumerate}

\end{frame}

%==============================================================================
% SLIDE 11: Conclusion
%==============================================================================
\begin{frame}{Conclusion}

\begin{center}
\Large
\textbf{Cognitive processing varies with context and social resources}
\end{center}

\vspace{1em}

\textbf{Key insights:}
\begin{itemize}
    \item High conflict constrains people to System 1 processing
    \item Recovery periods allow socially connected individuals to engage System 2
    \item Social connections matter because they enable the mental effort needed for analytical thinking
    \item This variation matters for humanitarian response
\end{itemize}

\vspace{1em}

\begin{center}
\large
Rather than assuming uniform decision-making,\\
our model captures how environmental pressures and social resources\\
shape cognitive availability
\end{center}

\end{frame}

%==============================================================================
% SLIDE 12: Configuration (Section 7)
%==============================================================================
\begin{frame}{Configuration: Nuclear Case Parameters}

\begin{table}[h]
\small
\begin{tabular}{llcll}
\toprule
Parameter & Symbol & Domain & Role & Description \\
\midrule
\texttt{alpha} & $\alpha$ & $(0, \infty)$ & Cognitive (free) & Psi curve shape \\
\texttt{beta} & $\beta$ & $(0, \infty)$ & Cognitive (free) & Omega collapse rate \\
\texttt{kappa} & $\kappa$ & $(0, \infty)$ & Cognitive (free) & S2 sensitivity \\
\texttt{max\_move\_speed} & $v$ & $(0, \infty)$ &
Physical constraint (fixed) & Mean vehicle speed under evacuation congestion \\
\bottomrule
\end{tabular}
\caption{Nuclear case configuration. \texttt{max\_move\_speed} is a physical
constraint fixed at an empirically grounded value.}
\end{table}

\vspace{0.5em}
\paragraph{Physical constraint parameter.} \texttt{max\_move\_speed}
($v$, km per timestep) is not a free cognitive parameter. It represents
the mean vehicle travel speed under evacuation congestion and is fixed
at an empirically grounded value rather than calibrated. For the
Fukushima case, probe-car telematics data from Toyota Motor Corporation
indicate average vehicle speeds of 7--20~km/h during the evacuation
\citep{Shimazaki2012}, implying $v \approx 14$--$40$ km per 2-hour
timestep. The data-preferred value is identified from the calibration
grid search and fixed for all subsequent analyses. Its sensitivity is
reported in the supplementary material but it does not enter the
count of free cognitive parameters ($\alpha$, $\beta$, $\kappa$).

\end{frame}

%==============================================================================
% SLIDE 13: Thank You
%==============================================================================
\begin{frame}
\begin{center}
\Huge Thank You

\vspace{1em}

\Large Questions?

\vspace{1.5em}

\normalsize
Michael Puma\\
Columbia University Climate School\\
Center for Climate Systems Research
\end{center}
\end{frame}

%==============================================================================
% References
%==============================================================================
\begin{frame}[allowframebreaks]{References}
\bibliographystyle{plainnat}
\begin{thebibliography}{99}
\bibitem[Shimazaki et~al., 2012]{Shimazaki2012}
Shimazaki, T., et~al. (2012).
\newblock Cluster analysis and discriminant analysis for determining
post-earthquake road recovery patterns.
\newblock [Primary source --- to be verified against Toyota G-BOOK
telematics study of Fukushima road recovery].
\end{thebibliography}
\end{frame}

\end{document}

